\noindent
This section introduces the motivation for the project, defines the problem being addressed, and sets out the measurable objectives used to evaluate whether the work was successful. It also summarises the main contributions and the practical scope limitations of the final prototype.

\subsection{Background and Motivation} \label{sec:introduction_background}

Robotic manipulators are increasingly used in domains such as industrial automation, teleoperation, and robot-assisted surgery, where safe and effective interaction with physical environments is essential.
Although visual feedback provides important global information about the workspace, it alone may be insufficient for developing fine motor control, accurate force regulation, and intuitive understanding of contact conditions.

Haptic interfaces address this limitation by enabling bidirectional physical interaction between a user and a virtual or remote environment. By rendering forces, such systems can communicate contact, constraint, and interaction information that is difficult to convey visually alone \cite{giriApplicationBasedReviewHaptics2021}. In robotic manipulation tasks, this is particularly valuable for training applications, where users must learn not only where to move, but also how to regulate force and respond to contact conditions. Prior work has shown that haptic feedback can improve task performance and reduce excessive applied force, particularly for less experienced users \cite{bergholzBenefitsHapticFeedback2023}.

Despite these benefits, access to haptic systems remains limited in research and educational settings. Commercial interfaces provide mature and high-performance force feedback, but are typically priced substantially above academic budgets and are often distributed through vendor-specific SDK integration pathways \cite{giriApplicationBasedReviewHaptics2021}. This can restrict accessibility and make it difficult for researchers and students to experiment with alternative control strategies, custom simulation environments, or modified hardware configurations. As a result, there is strong practical motivation for low-cost and modular haptic platforms that are suitable for prototyping, experimentation, and robotics training research.

At the same time, achieving stable and convincing haptic interaction is technically challenging. Real-time haptic rendering commonly requires update rates on the order of 1~kHz together with low latency and carefully designed control structure in order to render contact stably and responsively \cite{ruspiniHapticDisplayComplex1997}. The stability of haptic interaction is particularly sensitive to delay in the feedback loop, since communication and computation delays introduce phase lag and reduce the maximum achievable virtual stiffness \cite{adamsStableHapticInteraction1999}. Consequently, the design of a practical haptic system is not simply a matter of adding actuators to an input device: it requires coordinated consideration of mechanical design, embedded control, communication architecture, simulation coupling, and force rendering strategy.

\subsection{Problem Definition}

This dissertation addresses the design and evaluation of a low-cost real-time haptic interface system for robotic manipulation training, with explicit consideration of practical implementation constraints.
The specific challenge is not merely to construct a haptic input device, but to integrate multiple subsystems into a bounded bidirectional interaction loop: user motion must be sensed at the device, transmitted to a host system, interpreted within a simulated environment, converted into interaction forces, mapped back into actuator commands, and rendered to the user in real time. In practice, each stage of this loop introduces constraints. Limited actuator capability affects the achievable force output, communication delay limits responsiveness and achievable virtual stiffness, and software architecture determines whether the system remains extensible and maintainable.

Existing work demonstrates the value of haptic feedback and the feasibility of low-cost haptic hardware, but there remains a practical gap between simple hardware demonstrations and fully integrated systems that combine low-cost actuation, real-time communication, physics-based simulation, and explicit treatment of latency and stability. Educational and open-source devices have shown that affordable haptic hardware is possible \cite{martinez3DPrintedHaptic2016,martinezOpenSourceModular2017}, and low-cost 2-DOF academic systems have also been demonstrated \cite{lavatelliDesignOpenSource2014}. However, comparatively less emphasis is placed on complete end-to-end architectures in which hardware, embedded control, host-side simulation, rendering, and communication are developed together and evaluated as an integrated real-time system. Communication architecture is therefore a first-order design constraint in this project, since buffering, scheduling, and transport delay can directly affect both stability and contact fidelity.

The problem addressed in this project is therefore defined as follows: \emph{to design, implement, and evaluate a low-cost haptic rendering system that is sufficiently modular for experimentation, sufficiently responsive for real-time interaction, and capable of bounded, controllable bidirectional haptic coupling with a simulated robotic environment}.

\subsection{Aim and Scope of This Dissertation}

The project focuses on demonstrating the feasibility of a complete end-to-end haptic architecture, covering: device design and construction; embedded firmware development for sensing, communication, and actuator control; implementation of a host-side framework integrating physics simulation, haptic rendering, and visualisation; and experimental evaluation of communication performance, timing, bounded contact behaviour, and practical interaction quality.
The scope is intentionally limited. The dissertation does not attempt fully calibrated force output, commercial-level rendering fidelity, or a formal user study. Final hardware validation was performed with force feedback on one joint only, due to actuator failure during development. Consequently, conclusions about full 2-DOF force rendering are provisional, while the communication, host-side timing, simulation, and safety findings remain directly applicable because they do not depend on two-axis actuation. The principal contribution therefore lies in \emph{system-level design, integration, and evaluation}, together with analysis of the practical constraints that limit achievable performance. As the results demonstrate, the dominant limitation is communication latency rather than host-side computation, and addressing this is identified as the primary route to higher-fidelity interaction.

\subsection{Objectives}

To achieve this aim, the project pursued six SMART objectives:

\begin{enumerate}[label=(\roman*)]
	\item Build a low-cost physical haptic interface with a total hardware cost below £200, capable of measuring joint position and accepting actuator torque commands during closed-loop operation.
	\item Achieve reliable host--device data exchange at a target update rate of 1\,kHz, with packet integrity assessed using measured receive rate, packet loss, out-of-order events, and round-trip latency.
	\item Demonstrate that host-side computation is not the dominant timing bottleneck by measuring mean closed-loop host processing latency below 1\,ms during integrated operation.
	\item Validate the host-side rendered contact model quantitatively using logged rendered-force data, showing a near-linear force--penetration relationship with fitted stiffness within 5\% of the configured virtual coupling value and high goodness of fit.
	\item Verify safety behaviour by measuring force saturation, embedded torque limiting, torque slew-rate limiting, and watchdog torque removal during communication interruption.
	\item Evaluate physical interaction performance under motor-enabled contact and constrained-motion tests, reporting penetration bounds, rendered-force command consistency, and limitations affecting generalisation to full multi-axis force feedback.
\end{enumerate}

\subsection{Contributions of This Work}

Prior work on low-cost haptic platforms focuses primarily on hardware design~\cite{martinez3DPrintedHaptic2016,martinezOpenSourceModular2017,lavatelliDesignOpenSource2014} or on isolated software demonstrations, with comparatively little emphasis on fully integrated real-time systems in which hardware, firmware, communication, simulation, and rendering are co-developed and evaluated together. This work addresses that gap. Its principal contributions are:

\begin{itemize}
	\item \textbf{A complete end-to-end haptic rendering system.} Custom hardware, embedded firmware, a packet-based communication protocol, and a modular host-side framework with real-time physics and OpenGL visualisation are integrated into a single bidirectional pipeline and evaluated as a unified system. This type of co-design and cross-layer evaluation is less commonly emphasised in existing low-cost academic haptic prototypes.
	
	\item \textbf{Demonstration of bounded contact interaction under non-negligible communication delay.} Bounded, controllable contact behaviour was observed during the tested interactions despite a measured round-trip latency of approximately 15--16\,ms --- more than an order of magnitude above the 1\,ms control timestep. This behaviour was achieved without a formal passivity guarantee through the combined use of virtual coupling, damping, snapshot-based communication, multi-threaded decoupling, and an embedded torque slew-rate limiter.
	
	\item \textbf{Quantitative identification of transport delay, not computation, as the dominant performance constraint.} Instrumented runtime logging isolates host-side processing latency from transport-layer delay (as quantified in Section~\ref{sec:results_host_timing}), demonstrating that the communication interface, not the host processor, is the bottleneck. This finding directly motivates communication-architecture improvements as the primary route to higher haptic stiffness in this class of system.
	
	\item \textbf{Characterisation of USB CDC burst behaviour and a snapshot-based mitigation strategy.} Measured host-side packet inter-arrival times reveal strongly bimodal delivery: most packets arrive in near-simultaneous bursts separated by gaps of 14--16\,ms, consistent with USB CDC buffering. A snapshot-based channel design is shown to decouple the haptic rendering loop from this irregular delivery pattern, preventing stale-state accumulation without requiring changes to the underlying communication stack.
\end{itemize}

%\subsection{Structure of the Dissertation}

%The remainder of this dissertation is organised as follows.
%
%Chapter~2 reviews relevant literature on haptic interfaces, haptic rendering methods, stability considerations, and low-cost haptic hardware. Chapter~3 presents the overall system architecture. Chapter~4 describes the practical implementation of the hardware, firmware, communication protocol, and host-side software. Chapter~5 outlines the experimental methodology used for evaluation. Chapter~6 presents the results. Chapter~7 discusses the implications of these results, with particular focus on latency, stability, and hardware limitations. Finally, Chapter~8 concludes the dissertation and outlines opportunities for future work.
